% /root/monodromy, 2026-11-19. Root and RD complete proof audits passed; no uniform image claim.
\subsection{A local self-tacnode criterion for a nonordinary surface}
\label{subsec:ss-self-tacnode}
\status{Proved under the stated hypotheses. The parent and reduction/descent agent independently audited both complete local proofs. This subsection isolates
one local geometric step. It does not assert irreducibility of a Hasse divisor
or determine any new full monodromy image.}

\begin{lemma}[Two-point Cartier obstruction]
\label{lem:ss-two-point-Cartier}
Let $C$ be a smooth projective curve over an algebraically closed field
of characteristic two. Let $P\ne Q$ be points, with chosen uniformizers
$x_P,x_Q$, and put $A_C=\ker(C:H^0(\omega_C)\to H^0(\omega_C))$.
Let $U\subset H^0(\omega_C(2P+2Q))$ be the subspace of forms with
zero residues at both points. Taking the double-pole coefficients gives
\[
 0\longrightarrow H^0(\omega_C)\longrightarrow U
   \xrightarrow{\mathrm{pp}}k^2\longrightarrow0.
\]
Cartier maps $U$ to $H^0(\omega_C)$ and induces a semilinear map
\[
 \delta:k^2\longrightarrow\operatorname{coker}
       (C:H^0(\omega_C)\to H^0(\omega_C)).
\]
Under Serre and finite-Frobenius duality its dual, with the appropriate
scalar Frobenius twists, is the evaluation map
\[
          A_C\longrightarrow k^2,
          \qquad\eta\longmapsto(\eta(P),\eta(Q)).
\]
In particular a Cartier-killed form with equal, nonzero double-pole
coefficients exists if and only if
$\eta(P)+\eta(Q)=0$ for every $\eta\in A_C$.
\end{lemma}
\begin{proof}
Riemann--Roch gives $\dim H^0(\omega_C(2P+2Q))=g+3$.
The residue theorem imposes the one relation between the two residues,
and the residue map has rank one. Thus $\dim U=g+2$.
The only forms in $U$ with both double-pole coefficients zero are
holomorphic, proving the exact sequence. The local rule for Cartier
kills the double-pole term and takes the square root of the residue;
hence it has the asserted target.

For clarity we identify the dual boundary, rather than inferring it
from a dimension count. Finite Frobenius duality identifies the full
meromorphic Cartier map with
\[
 H^1(\mathcal O_C(-P-Q))\xrightarrow{F}
 H^1(\mathcal O_C(-2P-2Q)).
\]
The respective kernels of the maps to $H^1(\mathcal O_C)$ are the
constant-difference line $k^2/k$ and
$H^0(\mathcal O_{2P+2Q})/k$. Frobenius is bijective on the
constant-difference line, so it can be eliminated before taking the
remaining two first-jet coordinates.

Explicitly, represent a Frobenius-kernel class on $H^1(\mathcal O_C)$
by a Cech cocycle $c_{ij}$, with $c_{ij}^2=b_i-b_j$.
Then $db_i$ glue to a global Cartier-killed form $\eta$.
The sequence
$0\to\mathcal O_C\xrightarrow F F_*\mathcal O_C
 \xrightarrow d B_C^1\to0$
and bijectivity of Frobenius on constants identify all such classes
with $A_C$, with the scalar twist just indicated. Choose disjoint
small affine neighbourhoods of $P,Q$ in the cover. Lifting the class
to $H^1(\mathcal O_C(-P-Q))$ is ambiguous by the
constant-difference line. There is a unique adjustment of that line
which kills the constant difference of its Frobenius image.
Subtracting one common constant then makes both zeroth-order terms
of $b$ zero. The remaining boundary has first-jet coordinates
\[
                   db(P)=\eta(P),\qquad db(Q)=\eta(Q).
\]
Serre duality pairs the coefficients of $x^{-2}dx$ with these
first-jet coordinates, by the residue of their product. This gives
the stated evaluation map. This is the two-point version of the
fully written one-point calculation in Lemma~\ref{lem:pointed-Cartier};
the finite-duality identification is
\cite[Lemma~1.3.6 and Proposition~1.3.7, pp.~26--27]{BrionKumar2005}.
Pairing the vector $(1,1)$ with that evaluation map proves the last
assertion. All relevant assertions are vanishing assertions and are
unchanged by the scalar square roots in Cartier.
\end{proof}

\begin{lemma}[Smoothness at a mixed self-tacnode]
\label{lem:ss-mixed-self-tacnode}
Let $A',A$ be superspecial abelian surfaces in characteristic two,
and let $\phi:A'\to A$ be an odd-degree isogeny. Fix invariant volume
forms and identify the invariant one-forms using $\phi^*$.
Let $C'\subset A'$ be a smooth ample curve and suppose its image
$Z=\phi(C')$ is birationally normalized by $C'$. Assume that $Z$ has
ordinary nodes and exactly one simple tacnode, the latter obtained
by identifying $P,Q=P+g$ for a nonzero $g\in\ker\phi$.
Their tangent lines agree after translation. Choose the same tangent
coordinate on these two branches.

Assume that
\[
 \ker C_{C'}=
       H^0(A',\Omega^1_{A'})|_{C'}\oplus k\eta,
       \qquad \eta(P)\ne\eta(Q),
\]
and that $\mathcal N=\mathcal O_A(Z)$ has a separable polarization.
Suppose also that $\mathcal N$ is very ample. Then the relative
coherent Frobenius on $|\mathcal N|$, on the quotient by ambient
$H^1(A,\mathcal O_A)$, has corank one at $Z$, and its determinant
has a nonzero first derivative there. Thus the relative Hasse divisor
is smooth at $Z$.
\end{lemma}
\begin{proof}
Invariant one-forms take equal values on the two translated tangent
vectors. The hypothesis on $\eta$ and
Lemma~\ref{lem:ss-two-point-Cartier} show that no
Cartier-killed meromorphic form on $C'$ has equal nonzero
principal coefficients at $P,Q$.

Here are the dualizing conditions at the singularities. At the
tacnode, choose completed coordinates in which the two branches
are $y=0$ and $y=\varphi(x)$, where
$\varphi(x)=c x^2+O(x^3)$ and $c\ne0$.
The completed local ring is the subring of
$k[[x]]\oplus k[[x]]$ consisting of pairs congruent modulo $x^2$.
The adjunction generator restricts on both branches to
$dx/\varphi(x)$. Thus a dualizing form has at most double poles,
with equal double-pole coefficients and equal residues.
At a node the two poles are simple, with opposite residues.
Cartier acts on the normalization componentwise by the local trace
rule and takes square roots of residues. Consequently a
Cartier-killed dualizing form has zero residues at every branch.
The two-point lemma now excludes a nonzero double-pole coefficient
at the tacnode. It follows that
\[
             \ker C_Z=\ker C_{C'},\qquad \dim\ker C_Z=3.
\]

We identify the ambient boundary precisely. For every effective
Cartier divisor $Z$ on an abelian surface the adjunction sequence is
\[
 0\longrightarrow\mathcal O_A\xrightarrow{s}\mathcal N
 \longrightarrow\omega_Z\longrightarrow0,
\]
using the chosen invariant volume form. Its boundary on the
restriction of invariant one-forms is, after contraction with that
volume form, the differential of the polarization
$\lambda_{\mathcal N}:A\to A^\vee$.
Indeed, on trivializing opens write $s_i=g_{ij}s_j$. For an invariant
vector field $v$, the local lift of its contracted one-form is
$v(s_i)$. The boundary of these lifts is
$v(g_{ij})/g_{ij}$, which is exactly the differential of
$a\mapsto t_a^*\mathcal N\otimes\mathcal N^{-1}$.
Separability of the polarization makes this two-dimensional map
an isomorphism. In particular the restrictions of the two invariant
forms are independent, as required in the displayed direct sum.

Ample vanishing gives $H^1(A,\mathcal N)=0$. Thus the image of
$H^0(\mathcal N)/ks$ in $H^0(\omega_Z)$ is the kernel of that
boundary. It is Cartier-stable, and its Cartier kernel is
one-dimensional: subtract a unique invariant form from $\eta$
to put it in the boundary kernel. Denote the resulting form by
$\eta_0$. Its two values still differ, since invariant forms have
equal values. Lift $\eta_0$ to $u\in H^0(\mathcal N)$, uniquely
modulo $ks$.

We finish with the first derivative, keeping track of the quotient.
In a local frame of $\mathcal N$, holomorphicity on both normalized
branches implies
\[
        u(P)=0,\qquad
        u=a y+b x^2+\text{terms irrelevant to the branch values}.
\]
The two branch values of $\eta_0$ are $b/c$ and $a+b/c$.
Thus $a\ne0$. Since $s=y(y-\varphi(x))$ has quadratic term $y^2$
and no $xy$ term, the Cartier trace of $sv$ at the singular point
vanishes for every regular $v$.

In root-removed coordinates, Cartier is the bilinear map
$M(s):W\to W^{(2)}$, with $W=H^0(\mathcal N)$, and the relative
map is
\[
                   W/ks\longrightarrow W^{(2)}/ks^{(2)}.
\]
The Frobenius-twisted evaluation at the singular point annihilates
its image and is nonzero on its target quotient; it spans the
one-dimensional left kernel. Very ampleness supplies a global
$\dot s$ with zero value and first jet $x$ at that point. The
$xy$ coefficient of $\dot s\,u$ is $a\ne0$, and hence
\[
             \operatorname{ev}^{(2)}_P(M(\dot s)u)\ne0.
\]
This is the intrinsic derivative of the quotient map: a derivative
of the chosen lift $u$ contributes $M(s)\dot u$, annihilated by
evaluation; the target kernel line $ks^{(2)}$ has zero first
variation in characteristic two, and its value at the singular
point is zero. The rank-one adjugate formula therefore gives a
nonzero derivative of the relative determinant. The calculation
is on the affine coefficient cone; scaling contributes only the
radial derivative, which vanishes at a determinant-zero point, so
the conclusion descends to the projective system.
\end{proof}

\begin{remark}[The remaining global gate]
\label{rem:ss-self-tacnode-global-gate}
The lemma requires a specified self-tacnode with
$\eta(P)\ne\eta(Q)$. It does not prove that such points occur on
every translation-collision divisor of a norm family, and it does
not address its triple-identification divisors. Those incidence
statements remain unproved. In particular this local criterion
alone does not imply irreducibility of any higher-degree Hasse
divisor or a uniform supersingular full-group theorem.
\end{remark}
